% from ega3-1-fr.tex
\oldpage[0\textsubscript{III}]{25}
Rappelons que si $(F^p(X))_{p\in\mathbb Z}$ est une filtration (décroissante)
d'un objet $X\in\mathcal C$, on dit que cette filtration est \emph{séparée}
si $\inf(F^p(X))=0$, \emph{discrète} s'il existe un $p$ tel que $F^p(X)=0$,
\emph{exhaustive} (ou \emph{co-séparée}) si $\sup(F^p(X))=X$,
\emph{co-discrète} s'il existe $p$ tel que $F^p(X)=X$.

Nous dirons qu'une suite spectrale $E=(E_r^{pq},E^n)$ est \emph{faiblement
convergente} si l'on a
\[
  B_\infty(E_2^{pq})=\sup_k(B_k(E_2^{pq})),\qquad
  Z_\infty(E_2^{pq})=\inf_k(Z_k(E_2^{pq}))
\]
(autrement dit les objets $B_\infty(E_2^{pq})$ et $Z_\infty(E_2^{pq})$
sont déterminés par les données a) à c) de la suite spectrale E). Nous dirons
que la suite spectrale E est \emph{régulière} si elle est faiblement
convergente et si en outre~:
\begin{enumerate}
  \item[$1^\circ$] Pour tout couple $(p,q)$, la suite décroissante
    $(Z_k(E_2^{pq}))_{k\geq2}$ est \emph{stationnaire}~; l'hypothèse que E
    est faiblement convergente entraîne alors
    $Z_\infty(E_2^{pq})=Z_k(E_2^{pq})$ pour $k$ assez grand (dépendant de
    $p$ et $q$).
  \item[$2^\circ$] Pour tout $n$, la filtration
    $(F^p(E^n))_{p\in\mathbb Z}$ de $E^n$ est \emph{discrète} et
    \emph{exhaustive}.
\end{enumerate}

On dit que la suite spectrale E est \emph{co-régulière} si elle est
faiblement convergente et si en outre~:
\begin{enumerate}
  \item[$3^\circ$] Pour tout couple $(p,q)$, la suite croissante
    $(B_k(E_2^{pq}))_{k\geq2}$ est \emph{stationnaire}, ce qui entraîne
    $B_\infty(E_2^{pq})=B_k(E_2^{pq})$, et par suite
    $E_\infty^{pq}=\inf E_k^{pq}$.
  \item[$4^\circ$] Pour tout $n$, la filtration de $E^n$ est
    \emph{co-discrète}.
\end{enumerate}

Enfin, on dit que E est \emph{birégulière} si elle est à la fois régulière
et co-régulière, autrement dit si on a les conditions suivantes~:
\begin{enumerate}
  \item[(a)] Pour tout couple $(p,q)$, les suites
    $(B_k(E_2^{pq}))_{k\geq2}$ et $(Z_k(E_2^{pq}))_{k\geq2}$ sont
    \emph{stationnaires} et l'on a
    $B_\infty(E_2^{pq})=B_k(E_2^{pq})$ et
    $Z_\infty(E_2^{pq})=Z_k(E_2^{pq})$ pour $k$ assez grand (ce qui entraîne
    $E_\infty^{pq}=E_k^{pq}$).
  \item[(b)] Pour tout $n$, la filtration $(F^p(E^n))_{p\in\mathbb Z}$ est
    \emph{discrète} et \emph{co-discrète} (ce qu'on exprime aussi en disant
    qu'elle est \emph{finie}).
\end{enumerate}

Les suites spectrales définies dans (T, 2.4) sont donc les suites spectrales
birégulières.
\end{env}

\begin{env}[11.1.4]
\label{0.11.1.4-fr}
Supposons que dans la catégorie $\mathcal C$, les limites inductives filtrantes
existent et que le foncteur $\varinjlim$ soit \emph{exact} (ce qui équivaut à
dire que l'axiome AB 5) de (T, 1.5) est vérifié (cf. T, 1.8)). La condition
que la filtration $(F^p(X))_{p\in\mathbb Z}$ d'un objet $X\in\mathcal C$ est
exhaustive s'écrit alors aussi
$\varinjlim_{p\to-\infty}F^p(X)=X$. Si une suite spectrale E est faiblement
convergente, on a
$B_\infty(E_2^{pq})=\varinjlim_{k\to\infty}B_k(E_2^{pq})$~; si en outre
$u:E\to E'$ est un morphisme de E dans une suite spectrale $E'$ de
$\mathcal C$ faiblement convergente, on a
$u_2^{pq}(B_\infty(E_2^{pq}))=B_\infty(E_2^{\prime pq})$ en raison de
l'exactitude de $\varinjlim$. De plus~:
\end{env}

\begin{proposition}[11.1.5]
\label{0.11.1.5-fr}
Soient $\mathcal C$ une catégorie abélienne dans laquelle les limites
inductives filtrantes sont exactes, E, $E'$ deux suites spectrales régulières
de $\mathcal C$, $u:E\to E'$ un morphisme de suites spectrales. Si les
$u_2^{pq}$ sont des isomorphismes, il en est de même de $u$.
\end{proposition}

\begin{proof}
Nous savons déjà (\hyperref[0.11.1.2-fr]{11.1.2}) que les $u_r^{pq}$ sont des
isomorphismes et que
\[
  u_2^{pq}(B_\infty(E_2^{pq}))=B_\infty(E_2^{\prime pq})~;
\]
